\subsection{Discussion}

\subsubsection{Thrashing Mitigation: App-Aware UVM Eviction}
\label{uvm:sec:future:thrash}
As demonstrated in our evaluation,
    thrashing remains a primary bottleneck for UVM under heavy oversubscription.
A root cause of this thrashing is the semantic gap between
    the UVM driver's eviction policy and the framework's user-space caching allocator.
The driver typically evicts pages based on hardware-level access history or allocation order,
    which frequently conflicts with logical reuse:
    a block allocated via \texttt{cudaMallocManaged} long ago
    might be logically hot because the caching allocator just reassigned it to a new tensor.
To fully mitigate thrashing,
    the UVM driver needs eviction policies that are informed by
    application-level usage patterns and logical block lifetimes.

\subsubsection{Thrashing Mitigation: Utilization of App Hints}
\label{uvm:sec:future:hints}
Existing application-level offloading frameworks
    (e.g., for KV caches and model weights)
    already compute detailed schedules for when data will be needed next.
Rather than using these schedules to manually orchestrate data movement,
    frameworks could pass them down as hints to improve the UVM driver's
    eviction and prefetching policies.
While relying on application hints may seem contrary to UVM's goal
    of transparent memory management,
    it strikes a pragmatic balance:
    the application provides high-level guidance on data lifetimes,
    and the driver retains the responsibility of executing
    the physical migrations and managing the hardware page tables.


\subsubsection{Multiple GPU and Cross-Node Support}
Existing \uvmname is based on single GPU attached to single CPU node.
In practice, multi-GPU system is a hard requirement for AI workloads in production.
We need to test UVM and extend \uvmname to support multi-GPU system.
Existing Pytorch manages memory for each GPU separately;
in other words, each GPU has its own caching allocator that is unaware of other GPUs.

For UVM, the pattern will change since every GPU can access the memory. 


\subsubsection{Using CUDA Graphs to Hint UVM Migration}
UVM suffers from a lack of application-level information:
    the driver reacts to page faults after they occur,
    with no advance knowledge of which pages the next kernel will touch.
CUDA Graphs offer a natural way to close this gap.
Rather than launching kernels one at a time,
    a CUDA Graph captures an entire sequence of kernels ahead of time
    and replays the whole graph on each invocation;
    since DNN training and inference are highly repetitive,
    the access pattern of one replay closely predicts the next.
We can exploit this with a record-and-replay approach to migration:
    during the first execution of a captured graph,
    the runtime records the page faults and resulting migrations
    to build a per-graph access trace,
    and on subsequent replays the driver consults this trace to
    prefetch pages before they fault and to evict pages that are
    dead rather than those reused later in the graph.
This supplies the driver with the logical-lifetime information called for in
    Sections~\ref{uvm:sec:future:thrash} and~\ref{uvm:sec:future:hints},
    but derived automatically from graph capture rather than
    explicit application annotations.

\subsubsection{Concluding Remarks}
\label{uvm:sec:conclusion}
LLM inference is increasingly bound by GPU memory,
    and today each framework answers this with its own hand-tuned offloading code.
UVM promises to retire that machinery with transparent oversubscription,
    but has long been dismissed as too slow for production.
{\uvmname} shows that this gap is not fundamental:
    by propagating application-level intent down to the driver
    through a UVM-aware caching allocator and framework-guided eviction,
    it turns UVM into a practical substrate for PyTorch,
    matching the native \texttt{cudaMalloc} allocator when memory fits
    and running to completion---often far faster than default UVM---%
    where the native allocator simply crashes.
A performance gap to hand-tuned offloading remains under heavy oversubscription,
    but our thrashing-mitigation study and the directions above
    chart a concrete path toward closing it,
    and we hope {\uvmname} encourages frameworks to treat
    driver-managed memory as a first-class citizen rather than a last resort.
